% critique.tex — f172r, vv. 801–834 : commentaire critique
% Articles textuels repris mot pour mot du commentaire antérieur de l'auteur (voir repartition.md).
\cc{815}{ποθοῦντας.}{L'erreur manifeste du scribe du \italique{Ravennas}, copiant ποθοῦντας à la place de πο(ι)οῦντας, m'est l'occasion -- ὁ καιρός ! -- d'une courte digression portant sur les différentes graphies ποι- ou πο- lorsque la lettre qui suit est un `ο' ou un `ω'\nc{Débat qui a fait déjà couler beaucoup d'encre et que je ne compte ni trancher, ni même exposer en détail. Le lecteur intéressé pourra se reporter aux pages 226-230 que lui consacre Willem J. W. Koster dans son livre \italique{Autour d'un manuscrit d'Aristophane écrit par Demetrius Triclinius} (1957).}.
Je m'en tiens aux leçons de l'épigraphie attique ainsi qu'à l'usage scolaire : savoir, écrire systématiquement ποι-. Toutefois, le \italique{Ravennas} adopte majoritairement la graphie πο- lorsque la syllabe est brève. Cette convention, reprise par la plupart des éditions actuelles des \italique{Thesmophories}, est celle de nombreux manuscrits remontant à la période hellénistique. La confusion du scribe -- qui aurait lu rapidement une suite de `Ο' en y ajoutant un `Θ' -- tend à faire penser que cette règle était aussi de mise pour le manuscrit qu'il recopiait.}
\cc{819–23}{}{Le manuscrit présente le texte suivant :\citComm{καὶ μὲν δήπου καὶ τὰ πατρῷά γε\\
χείρους ἡμῶν εἰσιν σῴζειν·\\
ἡμῖν μὲν γὰρ σῶν ἔτι καὶ νῦν\\
τἀντίον ὁ κανὼν οἱ καλαθίσκοι\\
τὸ σκιάδειον·}Suit une liste de ces biens que les femmes sauvent mais que leur mari, eux, perdent. C'est là que le bât blesse.
D'abord, cette liste est étonnamment hétéroclite ; si les trois premiers éléments qu'elle comporte sont en rapport avec le filage et le tissage : τὸ ἀντίον, \italique{l'ensouple} (cylindre du métier à tisser) ; ὁ κανὼν, \italique{barre}, \italique{règle} (avec le sens spécialisé de \italique{tige de quenouille}), οἱ καλαθίσκοι, \italique{petits paniers} (où les femmes déposaient la laine filée), l'on se demande par quel ressort \italique{l'ombrelle}, τὸ σκιάδειον, vient clôturer cette liste.
Plus important : sur ces quatre éléments, seuls deux, κανών et σκιάδειον, subsistent dans la comparaison avec ce qu'en font les hommes. Pourquoi Aristophane nous parle-t-il alors de l'ἀντίον et des καλαθίσκοι, inutiles -- donc nuisibles -- si ce n'est pas pour filer la métaphore quelques vers plus loin ? (Je ne parle que de ce qui achoppe au niveau du sens. Quant à la forme, l'on avait déjà relevé l'inhabituelle succession de quatre brèves au vers 822.)\\
Aristophane veut nous peindre la sotte dissipation des hommes ; interrogeons les parémiographes\nc{Leutsch \& Schneidewin, \italique{Corpus Paroemiographorum Graecorum}, vol. I, Zenobius, II, 31. Recensé par Érasme (\italique{Adages}, n° 1088). Kassel \& Austin l'ont intégré comme \nieme{15} fragment de Platon le Comique. Suidas (α 4465 ; α 4466 ; α 4540) a relevé des expressions analogues : αὐταῖς ἀμάλαις ; αὐταῖσι διαβολαῖς ; αὐτοῖς ὅπλοις ; αὐτοῖς ταλάροις.} pour savoir comment on évoque une prodigalité de cigale ou une gestion de fourmi dans sa langue imagée et proverbiale :\citComm{Αὐτῷ κανῷ: αὕτη τάττεται κατὰ τῶν ἄρδην ὁτιοῦν λυμαινομένων. Φησὶ γὰρ ὁ Δίδυμος ὅτι τὰ δεῖπνα ἐπὶ κανῶν ἐκομίζετο. Καθάπερ ἐν Γρυψὶ λέγει καὶ Πλάτων·\\
Ἀνήρπακας ἁπαξάπαντ᾿ \textbf{αὐτῷ κανῷ}.\\
Et la corbeille elle-même : ce [proverbe] se range contre ceux prompts à se ruiner de quelque manière. Didyme dit que les repas étaient emportés dans des paniers. Ainsi Platon dans ses \italique{Griffons} :\\
Tu as tout emporté, tout, \italique{et le panier avec} !..}Le grec ajoute «~et le panier~» pour désigner tous les biens jusqu'au dernier\nc{Biens dévorés, en l'occurrence. Mais l'image du panier comme signe d'abondance se retrouve dans un autre proverbe, transmis par Hésychius (ε 7077) : εὔπλουτον κανοῦν.}, comme en français on ajoute «~et son père~» pour désigner «~tout le monde~». Dans les \italique{Grenouilles} (vers 560), ce n'est pas autrement qu'il faut entendre le désarroi de l'hôtesse face à Hercule, cet ogre qui a liquidé toutes ses provisions :\citComm{Μὰ Δί᾽ οὐδὲ τὸν τυρόν γε τὸν χλωρόν, τάλαν,\\
ὃν οὗτος \textbf{αὐτοῖς τοῖς ταλάροις} κατήσθιεν.\\
Par Zeus, et ce fromage frais, pauvre amie, qu'il dévorait \italique{avec les éclisses elles-mêmes}.}C'est peut-être encore ce même proverbe qu'on retrouve au vers 170 des \italique{Guêpes} dans la bouche de Philocléon :\citComm{μὰ τὸν Δί᾽ οὐ δῆτ᾽, ἀλλ᾽ ἀποδόσθαι βούλομαι\\
τὸν ὄνον ἄγων \textbf{αὐτοῖσι τοῖς κανθηλίοις}·\\
νουμηνία γάρ ἐστιν.\\
Par Zeus, certes non ! Je veux aller vendre l'âne \italique{avec ses paniers} ; car c'est nouvelle lune.}L'«~αὐτῇ λόγχῃ~» du vers 826 -- noter la présence du pronom --, mis comme permutant métaphorique d'«~αὐτῷ κανόνι~», rappelle le détournement du proverbe avec jeu de mots sur κανών / κανοῦν.\\
Ma conclusion est que le texte original des \italique{Thesmophories} -- l'\allemand{Urtext} -- comportait certainement au vers 822 le terme κανοῦν, \italique{corbeille}, et non κανών. La singulière présence des καλαθίσκοι, \italique{les corbeilles à laine}, s'explique alors par l'insertion fort importune dans le texte d'une glose de κανοῦν. \\
À la place des vers 822-823, je propose donc de lire :\citComm{καὶ μὲν δήπου καὶ τὰ πατρῷά γε\\
χείρους ἡμῶν εἰσιν σῴζειν·\\
ἡμῖν μὲν γὰρ σῶν ἔτι καὶ νῦν\\
ταὐτὸ κανοῦν\nc{Au vers 164, il y a confusion de finale -ων / -ουν du même genre entre les leçons du \italique{Ravennas} (διεκίνων) et de la \italique{Souda} ε 989 (διεκίνουν).} [καὶ] τὸ σκιάδειον·\\
Et assurément aussi, ils sont moins aptes que nous\\
à garder les traditions\\
qui sont, par notre fait, sauves encore aujourd'hui\\
-- même le panier [et] l'ombrelle.}}
